\section{主观体验的欠决定性与条件性桥接}
\label{sec:boundary}
系统在指定的干预下保留了一个比特，这确立了它的一项性质。能否由这些数据推出，这个具体系统在这个时刻具有主观体验？

\paragraph{T5.} 设操作理论 $T$ 的全部公理都不含体验谓词 $E(a,t)$。取该理论的任一模型 $M$，并指定一点 $u=(a,t)$。作两次扩张：第一次将 $E$ 解释为空集，第二次将其解释为一个包含 $u$ 的集合。因为改变的只有新增谓词，两种结构都满足 $T$，且所有操作性干预的结果都相同。在第一个结构中 $E(u)$ 为假，在第二个结构中为真。故即使修正规则的功能判据成立，$T$ 仍无法确定 $E(u)$ 的真值。这是关于理论表达能力的结论。构造出两个结构，尚不足以说明在同一载体上物理地实现两种不同体验是否可能。

\paragraph{T6.} 桥接公设组 $B=(B_{\mathrm{id}},B_{\mathrm{nom}},B_{\mathrm{epist}})$ 的状态为 \texttt{DECLARED\_NOT\_DERIVED}。三部分分别规定推断所指的载体、附加前提以及解释证据的规则。
\begin{description}[leftmargin=1.35cm,style=nextline]
\item[$B_{\mathrm{id}}$] 指定系统边界、宏观状态以及一组构成群的容许坐标重命名。描述在这些重命名之下保持不变。固定外部输入后，状态必须沿内部因果链影响后继状态。一条连续的历史谱系要求后继唯一。若对 KV 缓存作对称复制，随后两份副本独立演化，则记录一次 \texttt{BranchEvent}。两个因果后继拥有共同的前身，各自获得不同的分支标识。这是追踪历史的约定。体验主体在数值上是否同一，仍是开放问题。
\item[$B_{\mathrm{nom}}$] 记 $C=I_{\mathrm{cont}}\land C_{\mathrm{med}}\land A_{\mathrm{ind}}\land D_{\mathrm{dyn}}$。其中 $I_{\mathrm{cont}}$ 表示连续性。$C_{\mathrm{med}}$ 要求在屏蔽先前选择和信号所在的位置后仍存在内部中介，并能把同一规则迁移到另一类任务。$A_{\mathrm{ind}}$ 要求有独立检验表明多个模块能够访问同一内容，以及对可靠性的评估能够改变行动。$D_{\mathrm{dyn}}$ 要求处理过程引起持久的重组，并使这种重组在本轮结束后对后续处理产生因果影响。我们将充分条件 $\forall u\,[C(u)\Rightarrow E(u)]$ 作为假设。
\item[$B_{\mathrm{epist}}$] 状态 \texttt{UN} 表示现有证据仍使问题保持开放。决策标记 \texttt{NONE}、\texttt{INVESTIGATION\_PRIORITY} 和 \texttt{SENTIENCE\_CANDIDATE} 各自需要单独说明理由。要给出数值置信度，还需要先验概率和观测模型。
\end{description}
\begin{theorem}[条件定理 T6]
\label{thm:conditional-bridge}
在接受 $B_{\mathrm{nom}}$ 的前提下，由 $C(u)$ 可推出 $E(u)$。如果至少存在一个与证据和 $B_{\mathrm{id}}$ 相容的操作模型 $M$，且 $C(u)$ 在其中为假，那么加入 $B_{\mathrm{nom}}$ 后，同一组证据仍与 $E(u)$ 的两种真值相容。
\end{theorem}
\begin{proof}
第一项断言是在 $u$ 点应用蕴涵式。对第二项，在 $M$ 中取 $S=\{v:C(v)\}$，分别以 $E_0=S$ 和 $E_1=S\cup\{u\}$ 为解释，扩张同一个操作模型。两个解释都包含 $S$，因而都满足桥接公设。由于 $u\notin S$，二者在 $E(u)$ 的真值上不同。两种扩张的操作数据以及全部指定干预的结果都相同。
\end{proof}
缺少对某个合取项的测量，尚不能证明该项为假。在本研究中，对于规则~110、单比特寄存器、查表解码器和冻结的 OLMo，$E(a,t)$ 的操作性状态均为 \texttt{UN}。通过外部词元传递规则，仍留下内部中介及独立架构指标有待检验。计划中的 AI.04 检验可采用 Geiger 等人的因果抽象框架~\cite{geiger2025}。分布式对齐搜索（Distributed Alignment Search，DAS）为变量 $p$ 寻找候选子空间。在独立样例上，屏蔽原始位置 $h,e$，将供体历史的值代入这个子空间，就能检查回答是否遵循供体的规则。一轮计算中 KV 缓存的变化也不足以确立持久可塑性。
